\documentclass[letterpaper,11pt]{article}

% ==================================================
% PACKAGES
% ==================================================
\usepackage[utf8]{inputenc}
\usepackage[T1]{fontenc}
\usepackage[brazil]{babel}

\usepackage[left=1.5cm,right=1.5cm,top=1.5cm,bottom=1.5cm]{geometry}

\usepackage{titlesec}
\usepackage{enumitem}
\usepackage[hidelinks]{hyperref}
\usepackage{graphicx}
\usepackage{xcolor}

% ==================================================
% COLORS
% ==================================================
\definecolor{cvblue}{RGB}{20,40,90}
\definecolor{cvgray}{RGB}{90,90,90}

% ==================================================
% STYLE
% ==================================================
\pagestyle{empty}
\setlength{\parindent}{0pt}
\setlength{\parskip}{0pt}

\titleformat{\section}
{\large\bfseries\color{cvblue}}
{}{0em}{}
[\color{cvblue}\titlerule]

\setlist[itemize]{
leftmargin=1.2em,
itemsep=0pt,
topsep=1pt,
parsep=0pt,
partopsep=0pt,
label=\textcolor{cvblue}{\small$\bullet$}
}

% ==================================================
% MODULAR COMMANDS
% ==================================================

% ---------- header ----------
\newcommand{\cvheader}[5]{
\noindent
\begin{minipage}[t]{0.12\textwidth}
\vspace{0pt}
\includegraphics[width=2cm]{#1}
\end{minipage}
\hfill
\begin{minipage}[t]{0.85\textwidth}
\vspace{0pt}
{\LARGE \textbf{\color{cvblue}#2}}\\[1pt]
{\small #3}\\
{\small #4}\\
{\small #5}
\end{minipage}

\vspace{2pt}
{\color{cvblue}\rule{\textwidth}{0.6pt}}
\vspace{-10pt}
}

% ---------- section ----------
\newcommand{\cvsection}[2][6pt]{
\vspace{#1}
{\large\bfseries\color{cvblue} #2}\\[-6pt]
{\color{cvblue}\rule{\textwidth}{0.4pt}}
\vspace{-6pt}
}

% ---------- title line ----------
\newcommand{\cvtitle}[2]{
\textbf{#1} \hfill {\color{cvgray}#2}
}

% ---------- subtitle line ----------
\newcommand{\cvsubtitle}[2]{
\textit{\color{cvgray}#1} \hfill \textit{\color{cvgray}#2}
}

% ---------- institution block ----------
\newcommand{\cventry}[4]{
\cvtitle{#1}{#2}\\
\cvsubtitle{#3}{#4}
\vspace{2pt}
}

% ---------- paragraph ----------
\newcommand{\cvparagraph}[1]{
\vspace{2pt}
#1

\vspace{3pt}
}

% ---------- bullet item ----------
\newcommand{\cvbullet}[1]{
\item #1
}

% ---------- bullet list ----------
\newcommand{\cvliststart}{
\begin{itemize}
}

\newcommand{\cvlistend}{
\end{itemize}
}

% ---------- skill line ----------
\newcommand{\cvskill}[2]{
\textbf{\color{cvblue}#1:} #2\par
}

% ---------- language line ----------
\newcommand{\cvlanguage}[2]{
\textbf{\color{cvblue}#1:} #2\par
}

% ==================================================
% DOCUMENT
% ==================================================
\begin{document}

% ==================================================
% HEADER
% ==================================================
\cvheader
{logo.pdf}
{Gabriel dos Santos Schmitz}
{Backend Engineer (Go)\ |\ Distributed Systems\ |\ Applied Cryptography}
{(+55) 45 99936-2034\ |\ \href{mailto:gabrielzschmitz@protonmail.com}{gabrielzschmitz@protonmail.com}}
{\href{https://linkedin.com/in/gabrielzschmitz}{linkedin.com/in/gabrielzschmitz}\ |\ \href{https://github.com/gabrielzschmitz}{github.com/gabrielzschmitz}}

% ==================================================
% ABOUT
% ==================================================
\cvsection[-2pt]{Sobre}

\cvparagraph{
Desenvolvedor Backend com foco em Golang, sistemas distribuídos e engenharia de software orientada a performance.
Experiência profissional em desenvolvimento de aplicações SaaS, testes automatizados, arquitetura modular e pipelines CI/CD.
}

\cvparagraph{
Atuação complementar em sistemas low-level, criptografia aplicada e pesquisa em segurança pós-quântica,
incluindo implementação de primitivas criptográficas, blockchain e integração com ecossistemas como Hyperledger Fabric.
}

\textbf{\color{cvblue}Destaques}

\cvliststart
\cvbullet{Experiência profissional com backend em Golang aplicado a produtos SaaS}
\cvbullet{Pesquisa aplicada em criptografia pós-quântica e integração com blockchain}
\cvbullet{Implementação de sistemas low-level em C/C++ com foco em performance}
\cvbullet{Projetos públicos open source com adoção real e visibilidade técnica}
\cvbullet{Experiência com testes automatizados, CI/CD e arquitetura distribuída}
\cvlistend

% ==================================================
% EXPERIENCE
% ==================================================
\cvsection{Experiência Profissional}

\cventry
{NeticShard}
{Abril 2024 -- Março 2025}
{Desenvolvedor de Software Junior}
{Remoto -- LATAM}

\cvliststart
\cvbullet{Desenvolvimento backend de plataforma SaaS para resposta a incidentes em TI utilizando Golang.}
\cvbullet{Implementação de testes unitários, integração e E2E para aumento de confiabilidade do produto.}
\cvbullet{Estruturação de pipelines CI/CD e fluxo Git em ambiente colaborativo internacional.}
\cvbullet{Aplicação de princípios de Domain-Driven Design e arquitetura modular.}
\cvlistend

% ==================================================
% RESEARCH
% ==================================================
\cvsection{Pesquisa Acadêmica}

\cventry
{Rede Nacional de Pesquisa (RNP)}
{Novembro 2024 -- Novembro 2025}
{Bolsa de Pesquisa Científica}
{Toledo, PR}

\cvliststart
\cvbullet{Pesquisa em integração de criptografia pós-quântica em blockchain,
com foco na análise de eficiência e otimização de tamanhos de assinatura para
segurança de longo prazo.}
\cvbullet{Integração de algoritmos de criptografia pós-quânticas na linguagem Go.}
\cvbullet{Integração da linguagem Go modificada no Hyperledger Fabric.}
\cvlistend

\cventry
{Programa Institucional de Bolsas de Iniciação em Desenvolvimento Tecnológico e Inovação (PIBITI)}
{Setembro 2024 -- Novembro 2024}
{Bolsa de Pesquisa Científica}
{Toledo, PR}

\cvliststart
\cvbullet{Pesquisa em criptografia pós-quântica aplicada a reticulados e sistemas resistentes a ataques quânticos.}
\cvbullet{Estudo e implementação de primitivas criptográficas com foco em segurança, eficiência e aplicabilidade prática.}
\cvbullet{Estudo de KEMTLS e protocolos TLS pós-quânticos.}
\cvlistend

\cventry
{Universidade Tecnológica Federal do Paraná (UTFPR)}
{2023 -- 2024}
{Iniciação Científica}
{Toledo, PR}

\cvliststart
\cvbullet{Pesquisa em criptografia pós-quântica com foco em primitivas baseadas em reticulados.}
\cvbullet{Implementação do criptossistema Goldreich-Goldwasser-Halevi (GGH) e análise experimental de propriedades computacionais.}
\cvlistend

% ==================================================
% PROJECTS
% ==================================================
\newpage
\cvsection{Projetos}

\cventry
{Tomato.C}
{2022 -- 2025}
{C, NCurses, Makefile, Shellscript}
{}

\cvliststart
\cvbullet{Aplicação terminal em C puro com mais de 400 estrelas públicas no GitHub e adoção internacional por desenvolvedores.}
\cvbullet{Interface TUI com animações ASCII, notificações e arquitetura modular baseada em ncurses.}
\cvbullet{Distribuição multiplataforma com suporte a Linux, MacOS, WSL e integração com dependências externas.}
\cvlistend

\cventry
{ProofVote}
{2026}
{Blockchain, Criptografia, C++ / Sistemas distribuídos}
{}

\cvliststart
\cvbullet{Sistema experimental de votação eletrônica baseado em blockchain com integridade verificável e rastreabilidade criptográfica.}
\cvbullet{Implementação de estrutura de blocos, hashing, validação de votos e encadeamento de registros.}
\cvbullet{Exploração prática de consenso distribuído, tolerância a falhas e arquitetura blockchain do zero.}
\cvlistend

\cventry
{BibInject}
{2026}
{Python, Parsing de texto, Automação acadêmica}
{}

\cvliststart
\cvbullet{Ferramenta para automação de inserção e manipulação de referências bibliográficas LaTeX em páginas HTML.}
\cvbullet{Desenvolvido com TDD, cobertura de testes superior a 90\%, incluindo testes unitários e E2E.}
\cvbullet{Pipeline completo de CI/CD com GitHub Actions, incluindo lint, type checking e validação automatizada.}
\cvlistend

\cventry
{Motrix}
{2026 -- Presente}
{C++, Raylib, ECS Architecture}
{}

\cvliststart
\cvbullet{Implementação própria de Entity Component System minimalista voltada a motores gráficos.}
\cvbullet{Separação explícita entre componentes, sistemas e entidades com foco em performance.}
\cvbullet{Base arquitetural para experimentos visuais e simulações em tempo real.}
\cvlistend

\cventry
{Pendulum Tracker}
{2024}
{Python, Flask, Dash, OpenCV}
{}

\cvliststart
\cvbullet{Aplicação full-stack para rastreamento experimental de pêndulo em tempo real via visão computacional.}
\cvbullet{Extração automática de posição angular e geração de gráficos dinâmicos de oscilação.}
\cvbullet{Integração entre OpenCV e dashboard web interativo para análise física experimental.}
\cvlistend

% ==================================================
% EDUCATION
% ==================================================
\cvsection{Formação Acadêmica}

\cventry
{Universidade Tecnológica Federal do Paraná}
{2022 -- 2027}
{Bacharelado em Engenharia de Computação}
{Toledo, PR}

\cventry
{Universidade Estadual do Oeste do Paraná}
{2020 -- 2022}
{Programa de Ensino de Línguas -- Inglês}
{Cascavel, PR}

\cventry
{Harvard University \& Princeton University}
{2024}
{CS50x: Introduction to Computer Science \ |\ Algorithms, Part I \& II}
{Remoto}

% ==================================================
% SKILLS
% ==================================================
\cvsection{Competências Técnicas}

\cvskill{Linguagens}{Go, C, C++, Python, JavaScript, Java, SQL}
\cvskill{Backend}{REST APIs, Concurrency, Distributed Systems, Testing}
\cvskill{Infraestrutura}{Docker, CI/CD, Linux, AWS, Git}
\cvskill{Frameworks}{Flask, Node.js, HTMX}
\cvskill{Especialidades}{Cryptography, Systems Programming, Software Architecture}

% ==================================================
% LANGUAGES
% ==================================================
\cvsection{Idiomas}

\cvlanguage{Português}{Nativo}
\cvlanguage{Inglês}{Avançado (experiência profissional)}

\end{document}
